\documentclass[12pt]{article}
\usepackage[utf8]{inputenc}
\usepackage[T1]{fontenc}
\usepackage[margin=1in]{geometry}
\usepackage{setspace}
\usepackage{xcolor}
\usepackage[colorlinks=true, urlcolor=blue!60!black, linkcolor=blue!60!black]{hyperref}
\setlength{\parindent}{0em}
\setlength{\parskip}{0.6em}
\renewcommand{\baselinestretch}{1.0}
\begin{document}
\begin{center}
{\LARGE \textit{Ex parte Young}, 209 U.S. 123 (1908)\par}
\vspace{0.5em}
{\large Case Brief --- Two-Page Edition\par}
\end{center}

\section{Memory Jogger}
Federal courts may enjoin state officers from enforcing unconstitutional statutes; such an officer is ``stripped of his official or representative character'' and cannot claim the State's immunity (\href{https://www.courtlistener.com/opinion/96819/ex-parte-young/}{CourtListener}).

\section{Detailed Case Facts}
Minnesota cut railroad rates and backed the cuts with ruinous fines and prison terms for railroad employees. Northern Pacific shareholders sued Attorney General Edward T. Young in federal court, alleging confiscatory rates and penalties that closed the courts to any test of the laws.

\section{Procedural History}
The federal court enjoined Young from enforcing the Commodity Rate Act. He defied the injunction by obtaining state-court mandamus the next day, was held in contempt and taken into custody, and sought habeas corpus in the Supreme Court.

\section{Judicial Votes}
Petition dismissed 8--1. Justice Peckham wrote for the Court; Justice Harlan dissented.

\section{Holding}
A federal court may enjoin a state officer with some connection to enforcing an unconstitutional statute, even through general state-law duties; the State cannot immunize him from the Constitution. Penalties so severe as to deter testing a rate law in court are facially unconstitutional.

\section{Analysis of Opinions}
\textbf{Majority:} An officer enforcing an unconstitutional law acts without valid state authority. Unlike in \textit{Fitts v. McGhee}, Young had an enforcement connection through his general duties, and awaiting prosecution was an inadequate remedy given the crushing penalties. \textbf{Harlan:} the officer-State distinction is a fiction that would ``practically obliterate the Eleventh Amendment''; state courts could have heard the railway's federal defenses.

\section{Examples: Future Applications}
\textbf{Same-Side:} a clinic sues the health commissioner who enforces a law banning a federally protected procedure. \textbf{Opposite-Side:} an order to pay withheld welfare benefits from the treasury is barred (\textit{Edelman v. Jordan}). \textbf{Fence-Sitter:} a law enforced only by private suits, as in \textit{Whole Woman's Health v. Jackson}.

\section{Critique}
The ``double fiction'' treats the same act as state action for the Fourteenth Amendment but not the Eleventh. Progressives once decried ``government by injunction,'' yet \textit{Young} became the backbone of civil rights litigation, and the Court has steadily narrowed it.

\section{Key Quotations}
\begin{quote}
``If the act which the state Attorney General seeks to enforce be a violation of the Federal Constitution, the officer, in proceeding under such enactment, comes into conflict with the superior authority of that Constitution, and he is, in that case, stripped of his official or representative character, and is subjected in his person to the consequences of his individual conduct. The State has no power to impart to him any immunity from responsibility to the supreme authority of the United States.''
\end{quote}
\noindent --- Justice Peckham, Opinion of the Court, 209 U.S. at 159--60.

\end{document}
